\RequirePackage{iftex}
% DVI output (latex, uplatex or platex, then dvipdfmx): a global class option gives every package the dvipdfmx driver.
\def\officeparserdriver{}\ifpdf\else\ifXeTeX\else\def\officeparserdriver{dvipdfmx,}\fi\fi
\documentclass[\officeparserdriver 11pt]{article}
\iftutex
  \usepackage{fontspec}
\else
  \usepackage[T1]{fontenc}
  \usepackage[utf8]{inputenc}
  \usepackage{lmodern}
\fi
\usepackage[paperwidth=595.28pt,paperheight=841.89pt,top=72pt,bottom=72pt,left=72pt,right=72pt]{geometry}
\usepackage{parskip}
\usepackage{array,longtable}
\usepackage[hyperfootnotes=false]{hyperref}
\hypersetup{colorlinks=true,
  linkcolor=blue,
  urlcolor=blue,
  citecolor=blue,
  pdfcreator={officeParser},
  pdftitle={Sheet1}}
\setcounter{secnumdepth}{-\maxdimen}
\title{Sheet1}
\date{}

\begin{document}

\section{Sheet1}

{\scriptsize\setlength{\tabcolsep}{2pt}
\begin{longtable}{|*{13}{>{\raggedright\arraybackslash}p{\dimexpr(\linewidth-26\tabcolsep-14\arrayrulewidth)/13\relax}|}}
\hline
\# Created: 11/28/2025 & 4:36:44 PM &  &  &  &  &  &  &  &  &  &  &  \\
\hline
\# Modified: 12/26/2025 & 11:08:56 AM &  &  &  &  &  &  &  &  &  &  &  \\
\hline
 &  &  &  &  &  &  &  &  &  &  &  &  \\
\hline
 &  &  &  &  &  &  &  &  &  &  &  &  \\
\hline
 &  &  &  &  &  &  &  &  &  &  &  &  \\
\hline
 &  &  &  &  &  &  &  &  &  &  &  &  \\
\hline
Demonstration of DOCX support in~calibre~ &  &  &  &  &  &  &  &  &  &  &  &  \\
\hline
This document~demonstrates~the ability of the~calibre~DOCX Input plugin to convert the various typographic features in a Microsoft Word (2007 and newer) document. Convert this document to a modern~ebook~format, such as AZW3 for Kindles or EPUB for other~ebook~readers, to see it in action.~ &  &  &  &  &  &  &  &  &  &  &  &  \\
\hline
There is support for images, tables, lists, footnotes, endnotes,~links,~dropcaps~and~various~types~of text and paragraph level formatting.~ &  &  &  &  &  &  &  &  &  &  &  &  \\
\hline
To see the DOCX conversion in action, simply add this file to~calibre~using the~“Add Books”~button and then click “Convert”.~~Set the output format in the top right corner of the conversion dialog to EPUB or AZW3 and click~“OK”.~ &  &  &  &  &  &  &  &  &  &  &  &  \\
\hline
One, with~a very long~line to~demonstrate~that the hanging indent for the list is working correctly~ &  &  &  &  &  &  &  &  &  &  &  &  \\
\hline
Two~ &  &  &  &  &  &  &  &  &  &  &  &  \\
\hline
One~ &  &  &  &  &  &  &  &  &  &  &  &  \\
\hline
Two~ &  &  &  &  &  &  &  &  &  &  &  &  \\
\hline
Three~ &  &  &  &  &  &  &  &  &  &  &  &  \\
\hline
Four with~a very long~line to~demonstrate~that the hanging indent for the list is working correctly.~ &  &  &  &  &  &  &  &  &  &  &  &  \\
\hline
Five~ &  &  &  &  &  &  &  &  &  &  &  &  \\
\hline
Six~ &  &  &  &  &  &  &  &  &  &  &  &  \\
\hline
One~ &  &  &  &  &  &  &  &  &  &  &  &  \\
\hline
Two~ &  &  &  &  &  &  &  &  &  &  &  &  \\
\hline
We now resume our normal programming~ &  &  &  &  &  &  &  &  &  &  &  &  \\
\hline
Four~ &  &  &  &  &  &  &  &  &  &  &  &  \\
\hline
 &  &  &  &  &  &  &  &  &  &  &  &  \\
\hline
 &  &  &  &  &  &  &  &  &  &  &  &  \\
\hline
Text Formatting~ &  &  &  &  &  &  &  &  &  &  &  &  \\
\hline
Inline formatting~ &  &  &  &  &  &  &  &  &  &  &  &  \\
\hline
Here, we~demonstrate~various types~of inline text formatting and the use of embedded fonts.~ &  &  &  &  &  &  &  &  &  &  &  &  \\
\hline
Here~is~some~bold,~italic,~bold-italic,~underlined~and~struck~out~~text. Then, we have a superscript~and a subscript. Now we see some~red,~green~and~blue~text. Some text with a~yellow highlight. Some text in a~box. Some text~in~ &  &  &  &  &  &  &  &  &  &  &  &  \\
\hline
inverse video &  &  &  &  &  &  &  &  &  &  &  &  \\
\hline
. &  &  &  &  &  &  &  &  &  &  &  &  \\
\hline
A paragraph with styled text:~subtle~emphasis~~followed~by~strong text~and~intense emphasis.~This paragraph uses document wide styles for styling rather than inline text properties as~demonstrated~in the~previous~paragraph~—~calibre~can handle both with equal ease.~ &  &  &  &  &  &  &  &  &  &  &  &  \\
\hline
Fun with fonts~ &  &  &  &  &  &  &  &  &  &  &  &  \\
\hline
This document has embedded the Ubuntu font family. The body text is in the Ubuntu typeface, here is~some text in the Ubuntu Mono typeface, notice how every letter has the same width, even~i~and m. Every embedded font will automatically be embedded in the output~ebook~during conversion.~~ &  &  &  &  &  &  &  &  &  &  &  &  \\
\hline
Paragraph level formatting~ &  &  &  &  &  &  &  &  &  &  &  &  \\
\hline
You can do crazy things with~paragraphs, if~the urge strikes you. For~instance~this paragraph is~right aligned~and has a right border. It has also been given a light gray background.~ &  &  &  &  &  &  &  &  &  &  &  &  \\
\hline
For~the lovers~of poetry amongst you, paragraphs with hanging indents, like this often\hfil\break{}come in handy. You can use hanging indents to ensure that a line of poetry~retains~its individual identity as a line even when the screen~is~ too~narrow to display it as a single line. Not only does this paragraph have a hanging indent,~it~is also has an extra top margin, setting it apart from the preceding paragraph.~ &  &  &  &  &  &  &  &  &  &  &  &  \\
\hline
 &  &  &  &  &  &  &  &  &  &  &  &  \\
\hline
 &  &  &  &  &  &  &  &  &  &  &  &  \\
\hline
Tables~ &  &  &  &  &  &  &  &  &  &  &  &  \\
\hline
ITEM~ & NEEDED~ &  &  &  &  &  &  &  &  &  &  &  \\
\hline
Books~ & 1~ &  &  &  &  &  &  &  &  &  &  &  \\
\hline
Pens~ & 3~ &  &  &  &  &  &  &  &  &  &  &  \\
\hline
Pencils~ & 2~ &  &  &  &  &  &  &  &  &  &  &  \\
\hline
Highlighter~ & 2 colors~ &  &  &  &  &  &  &  &  &  &  &  \\
\hline
Scissors~ & 1 pair~ &  &  &  &  &  &  &  &  &  &  &  \\
\hline
Tables in Word can vary from the extremely simple to the extremely complex.~calibre~tries to do~its~best when converting tables. While you may run into trouble with the occasional table,~the vast majority of~common cases should be converted very well, as~demonstrated~in this section.~Note that for~optimum~results, when creating tables in Word, you should set their widths using percentages, rather than absolute units.~~To the left of this paragraph is a floating two column table with a nice green border and header row.~ &  &  &  &  &  &  &  &  &  &  &  &  \\
\hline
Now~let’s~look at a fancier table—one with alternating row colors and partial borders. This table is stretched out to take 100\% of the available width.~ &  &  &  &  &  &  &  &  &  &  &  &  \\
\hline
City or Town~ & Point A~ & Point B~ & Point C~ & Point D~ & Point E~ &  &  &  &  &  &  &  \\
\hline
Point A~ & —~ &  &  &  &  &  &  &  &  &  &  &  \\
\hline
Point B~ & 87~ & —~ &  &  &  &  &  &  &  &  &  &  \\
\hline
Point C~ & 64~ & 56~ & —~ &  &  &  &  &  &  &  &  &  \\
\hline
Point D~ & 37~ & 32~ & 91~ & —~ &  &  &  &  &  &  &  &  \\
\hline
Point E~ & 93~ & 35~ & 54~ & 43~ & —~ &  &  &  &  &  &  &  \\
\hline
 &  &  &  &  &  &  &  &  &  &  &  &  \\
\hline
 &  &  &  &  &  &  &  &  &  &  &  &  \\
\hline
Next, we see a table with special formatting in various locations. Notice how the formatting for the header row and sub header rows is preserved.~ &  &  &  &  &  &  &  &  &  &  &  &  \\
\hline
College~ & New students~ & Graduating students~ & Change~ &  &  &  &  &  &  &  &  &  \\
\hline
 & Undergraduate~ &  &  &  &  &  &  &  &  &  &  &  \\
\hline
Cedar University~ & 110~ & 103~ & +7~ &  &  &  &  &  &  &  &  &  \\
\hline
Oak Institute~ & 202~ & 210~ & -8~ &  &  &  &  &  &  &  &  &  \\
\hline
 & Graduate~ &  &  &  &  &  &  &  &  &  &  &  \\
\hline
Cedar University~ & 24~ & 20~ & +4~ &  &  &  &  &  &  &  &  &  \\
\hline
Elm College~ & 43~ & 53~ & -10~ &  &  &  &  &  &  &  &  &  \\
\hline
Total~ & 998~ & 908~ & 90~ &  &  &  &  &  &  &  &  &  \\
\hline
Source:~Fictitious data, for illustration purposes only~ &  &  &  &  &  &  &  &  &  &  &  &  \\
\hline
Next, we have something a little more complex, a nested table,~i.e.~a table inside another table. Additionally, the inner table has some of its cells merged.~The table is displayed horizontally centered.~ &  &  &  &  &  &  &  &  &  &  &  &  \\
\hline
One~ & Two~ & To the left is a table inside a table, with some cells merged.~ &  &  &  &  &  &  &  &  &  &  \\
\hline
Three~ &  &  &  &  &  &  &  &  &  &  &  &  \\
\hline
Four~ &  &  &  &  &  &  &  &  &  &  &  &  \\
\hline
 &  &  &  &  &  &  &  &  &  &  &  &  \\
\hline
 &  &  &  &  &  &  &  &  &  &  &  &  \\
\hline
 &  &  &  &  &  &  &  &  &  &  &  &  \\
\hline
We~end~with a~fancy~calendar, note how~much of the original formatting is preserved.~Note that this table will only display correctly on~relatively wide~screens. In general, very wide tables or tables whose cells have fixed width requirements~don’t~fare well in~ebooks.~ &  &  &  &  &  &  &  &  &  &  &  &  \\
\hline
December 2007~ &  &  &  &  &  &  &  &  &  &  &  &  \\
\hline
Sun~ &  & Mon~ &  & Tue~ &  & Wed~ &  & Thu~ &  & Fri~ &  & Sat~ \\
\hline
 &  &  &  &  &  &  &  &  &  &  &  & 1~ \\
\hline
 &  &  &  &  &  &  &  &  &  &  &  &  \\
\hline
2~ &  & 3~ &  & 4~ &  & 5~ &  & 6~ &  & 7~ &  & 8~ \\
\hline
 &  &  &  &  &  &  &  &  &  &  &  &  \\
\hline
9~ &  & 10~ &  & 11~ &  & 12~ &  & 13~ &  & 14~ &  & 15~ \\
\hline
 &  &  &  &  &  &  &  &  &  &  &  &  \\
\hline
16~ &  & 17~ &  & 18~ &  & 19~ &  & 20~ &  & 21~ &  & 22~ \\
\hline
 &  &  &  &  &  &  &  &  &  &  &  &  \\
\hline
23~ &  & 24~ &  & 25~ &  & 26~ &  & 27~ &  & 28~ &  & 29~ \\
\hline
 &  &  &  &  &  &  &  &  &  &  &  &  \\
\hline
30~ &  & 31~ &  &  &  &  &  &  &  &  &  &  \\
\hline
 &  &  &  &  &  &  &  &  &  &  &  &  \\
\hline
 &  &  &  &  &  &  &  &  &  &  &  &  \\
\hline
Structural Elements~ &  &  &  &  &  &  &  &  &  &  &  &  \\
\hline
Miscellaneous structural elements you can add to your document, like footnotes, endnotes,~dropcaps~and the like.~~ &  &  &  &  &  &  &  &  &  &  &  &  \\
\hline
Footnotes \& Endnotes~ &  &  &  &  &  &  &  &  &  &  &  &  \\
\hline
Footnotes1~and endnotesi~are automatically recognized and both are converted to endnotes, with backlinks for maximum ease of use in~ebook~devices.~ &  &  &  &  &  &  &  &  &  &  &  &  \\
\hline
Dropcaps~ &  &  &  &  &  &  &  &  &  &  &  &  \\
\hline
D~ &  &  &  &  &  &  &  &  &  &  &  &  \\
\hline
rop~caps are used to emphasize the leading paragraph at the start of a section. In~Word~it is possible to specify how many lines of text a drop-cap should use.~ &  &  &  &  &  &  &  &  &  &  &  &  \\
\hline
Links~ &  &  &  &  &  &  &  &  &  &  &  &  \\
\hline
Two kinds of links are possible, those that refer to an external~website~and those that refer to~locations~inside the document itself. Both are supported by~calibre. For example, here is a link pointing to the~calibre download page.~Then we have a link that points back to the section on~paragraph level formatting~in this document.~ &  &  &  &  &  &  &  &  &  &  &  &  \\
\hline
 &  &  &  &  &  &  &  &  &  &  &  &  \\
\hline
 &  &  &  &  &  &  &  &  &  &  &  &  \\
\hline
Images~ &  &  &  &  &  &  &  &  &  &  &  &  \\
\hline
Centered images like this are useful for large pictures that should be a focus of attention.~~ &  &  &  &  &  &  &  &  &  &  &  &  \\
\hline
There is no analogous technology in~ebooks, so the conversion will usually end up placing the image either centered or floating close to the point in the text where it was~inserted, not necessarily where it appears on the page in Word.~ &  &  &  &  &  &  &  &  &  &  &  &  \\
\hline
 &  &  &  &  &  &  &  &  &  &  &  &  \\
\hline
 &  &  &  &  &  &  &  &  &  &  &  &  \\
\hline
Lists~ &  &  &  &  &  &  &  &  &  &  &  &  \\
\hline
All types of lists are supported by the conversion,~with the exception of~lists that use fancy~bullets,~these get converted to regular bullets.~ &  &  &  &  &  &  &  &  &  &  &  &  \\
\hline
Bulleted List~ &  &  &  &  &  &  &  &  &  &  &  &  \\
\hline
One~ &  &  &  &  &  &  &  &  &  &  &  &  \\
\hline
Two~ &  &  &  &  &  &  &  &  &  &  &  &  \\
\hline
Numbered List~ &  &  &  &  &  &  &  &  &  &  &  &  \\
\hline
1 & One, with~a very long~line to~demonstrate~that the hanging indent for the list is working correctly~ &  &  &  &  &  &  &  &  &  &  &  \\
\hline
2 & Two~ &  &  &  &  &  &  &  &  &  &  &  \\
\hline
Multi-level Lists~ &  &  &  &  &  &  &  &  &  &  &  &  \\
\hline
1 & One~ &  &  &  &  &  &  &  &  &  &  &  \\
\hline
1 & Two~ &  &  &  &  &  &  &  &  &  &  &  \\
\hline
1 & Three~ &  &  &  &  &  &  &  &  &  &  &  \\
\hline
2 & Four with~a very long~line to~demonstrate~that the hanging indent for the list is working correctly.~ &  &  &  &  &  &  &  &  &  &  &  \\
\hline
3 & Five~ &  &  &  &  &  &  &  &  &  &  &  \\
\hline
2 & Six~ &  &  &  &  &  &  &  &  &  &  &  \\
\hline
A~Multi-level~list with bullets:~ &  &  &  &  &  &  &  &  &  &  &  &  \\
\hline
One~ &  &  &  &  &  &  &  &  &  &  &  &  \\
\hline
Two~ &  &  &  &  &  &  &  &  &  &  &  &  \\
\hline
This bullet uses an image as the bullet item~ &  &  &  &  &  &  &  &  &  &  &  &  \\
\hline
Four~ &  &  &  &  &  &  &  &  &  &  &  &  \\
\hline
Five~ &  &  &  &  &  &  &  &  &  &  &  &  \\
\hline
Continued Lists~ &  &  &  &  &  &  &  &  &  &  &  &  \\
\hline
i & One~ &  &  &  &  &  &  &  &  &  &  &  \\
\hline
ii & Two~ &  &  &  &  &  &  &  &  &  &  &  \\
\hline
An interruption in our regularly scheduled listing, for this essential and~very relevant~public service announcement.~ &  &  &  &  &  &  &  &  &  &  &  &  \\
\hline
iii & We now resume our normal programming~ &  &  &  &  &  &  &  &  &  &  &  \\
\hline
iv & Four~ &  &  &  &  &  &  &  &  &  &  &  \\
\hline
 &  &  &  &  &  &  &  &  &  &  &  &  \\
\hline
 &  &  &  &  &  &  &  &  &  &  &  &  \\
\hline
Charts &  &  &  &  &  &  &  &  &  &  &  &  \\
\hline
 & Series 1 & Series 2 & Series 3 &  &  &  &  &  &  &  &  &  \\
\hline
Category 1 & 4.3 & 2.4 & 2 &  &  &  &  &  &  &  &  &  \\
\hline
Category 2 & 2.5 & 4.4000000000000004 & 2 &  &  &  &  &  &  &  &  &  \\
\hline
Category 3 & 3.5 & 1.8 & 3 &  &  &  &  &  &  &  &  &  \\
\hline
Category 4 & 4.5 & 2.8 & 5 &  &  &  &  &  &  &  &  &  \\
\hline
\end{longtable}
}

% Comment: # TestString: Hello from custom props

% Comment: # TestNumber: 42

% Comment: # TestBool: true

% Comment: # Sheet: Sheet1

% Comment: # Sheet: Sheet2

% Comment: # Sheet: Sheet3

% Comment: # Sheet: Sheet4

% Comment: # Sheet: Sheet5

% Comment: # Sheet: Sheet6

% Comment: # Sheet: Sheet7

% Comment: # Sheet: Sheet8

% Comment: # Sheet: Chart Sheet

\end{document}
